% Ортогональное разложение x = Px + Qx по подпространствам L и L'.
\begin{figure}[htbp]
\centering
\begin{tikzpicture}[scale=1.0, >={Stealth[length=2.2mm]}]

% подпространство L (прямая под углом 20 градусов)
\draw[blue!55!black, thick] (-1.6,{-1.6*tan(20)}) -- (4.4,{4.4*tan(20)});
\node[blue!55!black, font=\footnotesize, anchor=west] at (4.45,{4.4*tan(20)})
  {$L = \{Px\}$};

% ортогональное подпространство L' (под углом 110 градусов)
\draw[orange!75!black, thick] ({-1.3*cos(110)},{-1.3*sin(110)})
                           -- ({3.0*cos(110)},{3.0*sin(110)});
\node[orange!60!black, font=\footnotesize, anchor=south east]
  at ({2.9*cos(110)},{2.9*sin(110)}) {$L' = \{Qx\}$};

% отметка прямого угла в начале координат
\draw[gray] ({0.32*cos(20)},{0.32*sin(20)})
  -- ({0.32*cos(20)+0.32*cos(110)},{0.32*sin(20)+0.32*sin(110)})
  -- ({0.32*cos(110)},{0.32*sin(110)});

\fill[black] (0,0) circle (1.2pt);

% вектор x и его проекции
\coordinate (O) at (0,0);
\coordinate (Py) at ({3.0*cos(20)},{3.0*sin(20)});     % Px
\coordinate (Qz) at ({1.6*cos(110)},{1.6*sin(110)});   % Qx
\coordinate (X)  at ({3.0*cos(20)+1.6*cos(110)},{3.0*sin(20)+1.6*sin(110)});

\draw[black, very thick, ->] (O) -- (X)
  node[midway, above left, font=\footnotesize, xshift=-2pt] {$x$};
\draw[blue!55!black, very thick, ->] (O) -- (Py)
  node[midway, below right, font=\footnotesize, yshift=-2pt] {$y = Px$};
\draw[orange!75!black, very thick, ->] (O) -- (Qz)
  node[midway, left, font=\footnotesize, xshift=-1pt] {$z = Qx$};

% пунктиры до параллелограмма
\draw[gray, dashed] (Py) -- (X);
\draw[gray, dashed] (Qz) -- (X);

\end{tikzpicture}
\caption{Ортогональное разложение из леммы~\ref{lem:projector}. Проектор $P$
отображает на подпространство $L$, дополнительный проектор $Q = E - P$ — на
ортогональное ему $L'$, и всякий вектор единственным образом представляется
как $x = Px + Qx$. В методе Розена $L$ — линейная оболочка строк матрицы
активных ограничений, а $L'$ — её ядро, то есть направления, не выводящие с
текущей грани.}
\label{fig:projection-decomposition}
\end{figure}
